\subsection{Datasets and Baselines}

We evaluate L-GTA and baselines on four time series domains with different frequencies and scales. \textbf{Tourism} (quarterly; Q) is the Australian Tourism dataset (TourismSmall), comprising quarterly visitor numbers by purpose, state, and city type. \textbf{Wiki2} (daily; D) consists of daily Wikipedia page-view series from the HierarchicalData benchmark. \textbf{Labour} (monthly; M) contains Australian monthly labour series from the same hierarchical suite. \textbf{M3} provides competition series in yearly (Y) and monthly (M) frequencies. Together these cover seasonal and non-seasonal, short- and long-horizon settings.

We compare L-GTA to the following. \textbf{Direct transformation} applies the same augmentations (jitter, scaling, magnitude warp, drift, trend) in observation space with fixed strength $\sigma$; it serves as an upper bound on controllability and transformation fidelity, since no latent round-trip is involved. \textbf{TimeGAN}, \textbf{TimeVAE}, and \textbf{Diffusion-TS} are generative baselines that produce synthetic series from fitted models without a notion of transformation type or strength; they are included in downstream forecasting and synthesis-quality comparisons. For controllability and signature analysis we focus on L-GTA variants (including the full model with and without dynamic features) versus Direct, using the same five transformations and $\sigma \in \{0.5, 1.0, 2.0\}$ where applicable.

The component ablation evaluates eight L-GTA variants to isolate the effect of latent structure, equivariance, encoder design, and dynamic features. \textbf{A: Global, No Equiv} uses a global (time-pooled) latent representation and no equivariance regularisation. \textbf{B: Temporal, No Equiv} keeps the full temporal latent representation but still without equivariance. \textbf{C: Global + Equiv} uses the global latent with equivariant decoder training (equivariance weight 1.0). \textbf{D: Full L-GTA} is the full model: temporal latent, equivariance, and dynamic (time-varying) features. \textbf{E: Simple Enc (temporal)} uses a simplified encoder with temporal latent and no equivariance. \textbf{F: Full + channel attn} adds channel attention to the full encoder but without equivariance. \textbf{G: Full L-GTA, no dynamic} is the full model with equivariance but with dynamic features disabled; this is the variant used in the cross-dataset signature and controllability figures. \textbf{H: Channel attn, no dynamic} combines channel attention with no equivariance and no dynamic features. All variants share the same five transformations, $\sigma$ sweep, and evaluation protocol so that mean $\rho$, monotonic \%, reconstruction MSE, and fingerprint distance are comparable across the grid.

\subsection{Results and Discussion}

In this section, we interpret the empirical results in light of the research questions introduced in Section~\ref{subsec:experimental_setup}. Our discussion is organised around controllability of latent transformations (\textbf{Q1}), downstream forecasting utility and comparison with state-of-the-art augmentation methods (\textbf{Q2}), synthesis quality and practical cost (\textbf{Q3}), and the contribution of individual architectural components (\textbf{Q4}).

\paragraph{Controllability of latent transformations (Q1).}
We ask whether varying the magnitude of a transformation in latent space produces a predictable, interpretable change in the decoded series. For each transformation type (e.g., jitter, scaling, magnitude warp), we sweep over latent magnitude (e.g., direction and scale) and decode to obtain synthetic series; we compare these to the same transformations applied directly in observation space. We quantify controllability in two ways: \textbf{mean $\rho$} is the Spearman correlation between the latent magnitude and an observed transformation strength (e.g., scale or variance of the residual); higher $\rho$ means that ``turning the knob'' in latent space reliably changes the decoded effect. \textbf{Monotonic \%} is the proportion of (dataset, frequency) cells in which that relationship is monotonic (higher is better). We also compute \textbf{fingerprint distance} to the corresponding direct transformation: the Euclidean distance between a 3-D residual fingerprint (autocorrelation, linearity, amplitude dependence) of the decoded output and that of the direct transformation; low distance means the decoded series exhibit the same qualitative transformation signature as the direct one.

Across datasets and frequencies, the full L-GTA model achieves a high mean Spearman correlation (mean $\rho = 0.95$) and a high proportion of monotonic responses (90\%), only slightly below the direct transformation baseline (mean $\rho = 1.0$, 100\% monotonic). At the same time, L-GTA substantially reduces the fingerprint distance to the target direct transformation compared with weaker variants, indicating that the decoded series closely match the intended qualitative behaviour (e.g., jitter-like or scaling-like residual patterns) without requiring hand-tuned transformation parameters.

\paragraph{Downstream forecasting utility and comparison to SOTA (Q2).}
To assess whether synthetic data from L-GTA and baselines actually help forecasting, we train forecasters on synthetic (or original) data and evaluate on held-out real series. We use two regimes: \textbf{TSTR} (train on synthetic, test on real), which measures how well the synthetic distribution supports learning a predictor that generalises to real data, and \textbf{augmentation}, where the forecaster is trained on real data augmented with synthetic series, which measures the added value of synthetic data when real data are limited. For each (dataset, frequency) cell we generate 3-variant synthetic sets with each method (L-GTA, Direct, TimeVAE, Diffusion-TS, TimeGAN) and also report results using only original data; we then train the same forecasting model per setting and evaluate with \textbf{MASE} (mean absolute scaled error; lower is better). Table~\ref{tab:downstream-summary} aggregates mean and standard deviation of the combined MASE across both regimes, and how often each method attains the best performance in a cell.

L-GTA attains the lowest mean combined MASE (2.99) among all methods and is the only approach that ranks first in either TSTR or augmentation in every evaluated cell. In contrast, strong generative baselines such as TimeVAE, TimeGAN, and Diffusion-TS, as well as the direct transformation baseline, exhibit higher mean errors and win fewer cells, indicating that their synthetic data either underfit or over-regularise the forecasting task. These results show that latent-space augmentation with L-GTA preserves and sometimes improves downstream predictive performance relative to using only real data.

\begin{table}[h]
\centering
\caption{Aggregated downstream forecasting results (3-variant, no-dynamic configuration). Mean and standard deviation of the combined MASE across TSTR and augmentation, and the number of cells in which each method attains the best performance.}
\label{tab:downstream-summary}
\begin{tabular}{lcccc}
\toprule
\textbf{Method} & \textbf{Mean MASE} & \textbf{Std MASE} & \textbf{1st in TSTR} & \textbf{1st in downstream / either} \\
\midrule
L-GTA & 2.99 & 1.35 & 3/5 & 5/5 \\
TimeVAE & 4.06 & 1.88 & 1/5 & 1/5 \\
Original & 4.12 & 1.83 & 0/5 & 0/5 \\
Direct & 4.44 & 2.25 & 1/5 & 1/5 \\
Diffusion-TS & 5.33 & 2.50 & 0/5 & 1/5 \\
TimeGAN & 5.98 & 3.35 & 0/5 & 0/5 \\
\bottomrule
\end{tabular}
\end{table}

\paragraph{Synthesis quality and practical cost (Q3).}
Beyond forecasting accuracy, we assess fidelity and diversity of the generated series. Table~\ref{tab:quality-summary} reports cross-dataset averages of key synthesis-quality metrics in the 3-variant setting. \textbf{Coverage} measures the proportion of the real data support that the synthetic distribution covers (diversity); higher values indicate that synthetic samples span the data manifold well. \textbf{Cosine similarity} compares the first principal component of real vs.\ synthetic features (fidelity); values close to 1 indicate that synthetic series align with the dominant real-data structure. The three distance metrics---\textbf{Wasserstein}, \textbf{Fréchet}, and \textbf{MMD}---quantify how close the synthetic distribution is to the real one in different ways (lower is better): Wasserstein and Fréchet emphasise location and spread, while MMD (maximum mean discrepancy) captures differences in moments in a reproducing-kernel Hilbert space.

L-GTA achieves high coverage (0.97) and near-perfect cosine similarity (1.00), and all three fidelity distances are very low (Wasserstein 0.01, Fréchet 0.01, MMD 0.02). TimeVAE and Diffusion-TS also attain low distances and high cosine similarity, but with lower coverage (0.83 and 0.75 respectively), suggesting their synthetic data may under-cover the support of the real distribution. TimeGAN reaches high coverage (0.95) but exhibits much larger Wasserstein, Fréchet, and MMD (0.56, 0.54, 0.56), indicating a systematic mismatch between its generated distribution and the real one despite good diversity. Overall, L-GTA combines strong diversity (high coverage) with the best fidelity (lowest distances and highest cosine similarity) among the compared methods.

\begin{table}[h]
\centering
\caption{Synthesis-quality metrics and computational cost. Values are averaged over all (dataset, frequency) pairs. Coverage and cosine similarity: higher is better; Wasserstein, Fréchet, and MMD: lower is better. Computational time reports overhead relative to L-GTA (baseline = 0).}
\label{tab:quality-summary}
\begin{tabular}{lccccccc}
\toprule
\textbf{Method} & \textbf{Coverage} & \textbf{Cos.\ sim.} & \textbf{Wass.} & \textbf{Fr\'echet} & \textbf{MMD} & \textbf{Comp.\ Time} \\
\midrule
L-GTA (3var, $\sigma{=}2.0$) & 0.97 & 1.00 & 0.01 & 0.01 & 0.02 & 0$\times$ \\
Direct (3var, $\sigma{=}2.0$) & 0.84 & 0.98 & 0.52 & 0.51 & 0.57 & 0$\times$ \\
TimeVAE (3var) & 0.83 & 0.99 & 0.01 & 0.01 & 0.03 & 8.2$\times$ \\
Diffusion-TS (3var) & 0.75 & 1.00 & 0.01 & 0.01 & 0.04 & 1085$\times$ \\
TimeGAN (3var) & 0.95 & 0.99 & 0.56 & 0.54 & 0.56 & 5487$\times$ \\
\bottomrule
\end{tabular}
\end{table}

From a practical standpoint, L-GTA provides this quality at a modest computational cost. For instance, on the Tourism-Q dataset, generating a 3-variant synthetic dataset with L-GTA takes less than a second and uses under 400\,MB of memory, whereas diffusion-based and TimeGAN baselines require several minutes of wall-clock time and more than 1.7\,GB of memory. While direct transformations remain essentially cost-free, they require dataset-specific manual tuning and, as discussed above, lead to inferior fidelity and privacy metrics.

\paragraph{Contribution of architectural components (Q4).}
To identify which design choices matter most for controllability and reconstruction, we systematically ablate architectural components of L-GTA and evaluate each variant under the same protocol. Table~\ref{tab:ablation-summary} reports cross-dataset averages for the main variants. \textbf{Mean $\rho$} is the Spearman correlation between the latent transformation magnitude and the observed transformation strength (higher is better); it measures how consistently the model maps latent changes to interpretable effects. \textbf{Monotonic \%} is the proportion of (dataset, frequency) cells in which that relationship is monotonic (higher is better). \textbf{Recon.\ MSE} is the mean squared reconstruction error when encoding and decoding real series (lower is better); it reflects how well the autoencoder preserves information. \textbf{Mean FP dist} (fingerprint distance) measures how closely the decoded transformation matches the target direct transformation in terms of residual structure (lower is better). For each transformation type (e.g., jitter, scaling, magnitude warp), a three-dimensional fingerprint is computed from the residuals (synthetic minus original): lag-1 autocorrelation (smooth vs.\ i.i.d.\ behaviour), linearity (R² of residuals vs.\ time, high for trend-like transforms), and amplitude dependence (correlation of absolute residuals with absolute original values, high for multiplicative transforms). The fingerprint distance is the Euclidean distance between the variant's fingerprint and the direct transformation's fingerprint; a low value indicates that the variant produces the intended qualitative transformation behaviour (e.g., jitter-like or scaling-like residuals) rather than a different signature.

Variants with equivariance regularisation and a global latent representation (e.g., Variant C: Global + Equiv) achieve substantially higher mean correlation and lower fingerprint distance than non-equivariant baselines, confirming the role of equivariance in enforcing consistent transformation effects. The full L-GTA model with dynamic features (Variant D) and L-GTA without dynamics (Variant G) both maintain high controllability (mean $\rho \geq 0.85$, monotonicity $\geq 80\%$), with only a small trade-off in reconstruction MSE. Removing dynamic features slightly reduces controllability but preserves low fingerprint distance, suggesting that these features mainly help align transformations with time-varying covariates rather than being strictly necessary for stable augmentation. These results support the design choice of combining a structured latent space with equivariant training to obtain controllable and data-efficient time series augmentation.

\begin{table}[h]
\centering
\caption{Component ablation summary. Metrics are averaged over all (dataset, frequency) pairs where each variant appears. FP dist = fingerprint distance to the corresponding direct transformation.}
\label{tab:ablation-summary}
\begin{tabular}{llcccc}
\toprule
\textbf{Variant} & \textbf{Description} & \textbf{Mean $\rho$} & \textbf{Monotonic \%} & \textbf{Recon.\ MSE} & \textbf{Mean FP dist} \\
\midrule
A & Global, No Equiv & 0.78 & 80.0 & 3.74M & 0.88 \\
B & Temporal, No Equiv & -- & 55.0 & 3.93M & 0.85 \\
C & Global + Equiv & 0.95 & 90.0 & 2.96M & 0.61 \\
D & L-GTA, with dynamic & 0.95 & 90.0 & 4.51M & 0.49 \\
E & Simple Enc (temporal) & -- & 85.0 & 5.10M & 0.83 \\
F & Full + channel attn & -- & 80.0 & 3.12M & 0.90 \\
G & L-GTA & 0.85 & 80.0 & 3.07M & 0.49 \\
H & Channel attn, no dynamic & 0.45 & 65.0 & 6.77M & 0.65 \\
\bottomrule
\end{tabular}
\end{table}
